%! Piecewise \& Step Functions — Unit 1, Lesson 1.4 — Warm-Up (Answer Key)
\documentclass[12pt]{article}
\usepackage{algebra2-article}
\usepackage{algebra2-key}   % pulls in -boxes; do NOT also load -boxes

\begin{document}

\pageheader{Unit 1, Lesson 1.4}{Warm-Up --- Answer Key}
% NAMESTRIP: no \namedateperiod here — the student writes their name once, on the cover.

\vspace{0.10in}

\begin{notesbox}{Getting Ready: Skills We'll Use Today}
Three things you already know. Today a function can follow \emph{more than one} rule; these are
the tools you'll need --- work quickly, then check with a neighbour.

\vspace{6pt}
\textbf{1. Evaluate the rule.} For $g(x)=2x-1$, find each output.
\begin{center}
\renewcommand{\arraystretch}{1.3}
\begin{tabular}{|c|c|c|c|}
\hline
$x$ & $-2$ & $0$ & $3$ \\ \hline
$g(x)=2x-1$ & \ans{$-5$} & \ans{$-1$} & \ans{$5$} \\ \hline
\end{tabular}
\end{center}
And for $h(x)=5-x$: \quad $h(2)=$ \ans{$3$} \qquad $h(-4)=$ \ans{$9$}

\vspace{8pt}
\textbf{2. The V is two lines (Lesson 1.3).} The parent $y=|x|$ is two straight lines pasted at
the vertex. Fill in each rule.
\begin{center}
$y=|x|$ behaves like $y=$ \ans{$x$} when $x\ge0$, \; and like $y=$ \ans{$-x$} when $x<0$.
\end{center}
Check: $|-6|=$ \ans{$6$}, and $-(-6)=$ \ans{$6$}. \; Same answer? \ans{yes}

\vspace{8pt}
\textbf{3. Endpoints and inequalities.} A number line uses a \emph{closed} dot $\bullet$ to
\emph{include} a point and an \emph{open} dot $\circ$ to \emph{exclude} it.
\begin{itemize}[leftmargin=1.6em, itemsep=4pt]
  \item Is $x=1$ a solution of $x<1$? \ans{no} \quad Is $x=1$ a solution of $x\le1$?
        \ans{yes}
  \item To show $x\ge2$ on a number line, the dot at $2$ is \ans{closed} (open / closed).
        To show $x>2$, it is \ans{open}.
  \item Suppose one rule is used when $x<1$ and a \emph{different} rule when $x\ge1$. The input
        $x=1$ uses \ans{the second rule} (the first rule / the second rule / both).
\end{itemize}
\end{notesbox}

\end{document}
